% Statement and cited-input packet. Target proof is separately embargoed.
% Frozen II source SHA256: 74497bfe327fae25c9cdcb17334b3bf42a9793a9df7c0a4404d97568f2f760f0
% Target owner: cdp:st:C.2.RAT.4; tag: 01CG

\begin{proposition}[Every rational degree-two tacnode cover has positive Euler number]\hypertarget{tag-01CG}{}\label{cdp:st:C.2.RAT.4}


Retain the torus-fibre, rational-normalization and specified
ruled-boundary morphism of \ref{cdp:st:C.2.RAT.1}. Suppose $C$ is not
contracted, $H=h(C)_{\rm red}$ has generic transverse type
$A_3$, and the actual map $C\to H^\nu$ has degree two.
Then $e(S)\ge1$. No smoothness or local complete intersection
assumption on the image curve is required.
\end{proposition}

% Cited input cdp:st:C.2.RAT.1; statement or labelled formula only.
\begin{proposition}[Euler signs for two of the rational boundary classes]\hypertarget{tag-01CD}{}\label{cdp:st:C.2.RAT.1}


Let $f:Y\to B$ be a proper holomorphic map from a smooth compact
complex threefold with $b_2(Y)=0$, with connected fibres and smooth
general fibre a complex two-torus. Let $mS$ be a full fibre with
irreducible nonnormal reduced support. Write
$h=\eta\pi:Z\to T\to S$ for the minimal resolution of its
normalization and suppose that $Z$ is rational. Use the effective
divisor $K_Z+D=0$ of \SeriesRef{cdp:st:A.5.16}.

If a contraction to $\mathbf P^2$ sends $D$ to three concurrent
reduced lines, then $e(S)\ge1$. More precisely,
$e(S)=-1+N_0+e(\Gamma_{\rm red})\ge1+N_0$, with $N_0$ as in
(A.RT.12).

Consider also a specified morphism $\rho:Z\to\mathbb F_n$,
$n\ge1$, which is a sequence of point blowups in reverse and has
\[
 \rho_*D=2C_0+\sum_{j=1}^r a_jF_j,\qquad
 C_0^2=-n,\qquad
 a_j\ge1,\quad \sum_j a_j=n+2,                  \tag{C.RA.1}\hypertarget{tag-00WQ}{}\label{cdp:eq:C.RA.1}
\]
where the $F_j$ are distinct ruling fibres. Let $C$ be the
strict transform of $C_0$. Suppose $h(C)$ is a curve and the
generic transverse singularity there is a cusp $A_2$.
Then $r\le3$ and $e(S)\ge1$. If $h(C)$ is smooth, then
$r\le2$. These assertions allow singular normalization,
arbitrary vertical conductor multiplicities and a nonreduced
conductor-image scheme.
\end{proposition}

% Cited input cdp:st:C.2.RAT.3; statement or labelled formula only.
\begin{proposition}[Every uncontracted degree-one horizontal image has nonnegative Euler number]\hypertarget{tag-01CF}{}\label{cdp:st:C.2.RAT.3}


Retain the torus-fibre, rational-normalization and ruled-boundary
morphism of \ref{cdp:st:C.2.RAT.1}. Suppose that $C$ is not contracted and
that its actual map to the normalization of
$H=h(C)_{\rm red}$ has degree one. Let $r_{\rm dom}$ be the
number of original ruling fibres $F_j$ whose positive vertical
boundary contains a strict conductor component mapping onto
$H$. Then
\[
       0\le r_{\rm dom}\le2,\qquad
       e(S)\ge2-r_{\rm dom}\ge0.
\tag{C.RA.21}\hypertarget{tag-00X3}{}\label{cdp:eq:C.RA.21}
\]
There is no assumption on the generic transverse type, on the
local complete intersection property of $H$, or on the
nilpotent structure of the conductor image.
\end{proposition}

